% !TeX root = ../main.tex
\chapter{CƠ SỞ LÝ THUYẾT VÀ NGHIÊN CỨU LIÊN QUAN}
\label{chap:cosolythuyet}

Chương này trình bày nền tảng lý thuyết của các thành phần được sử dụng trong quy
trình đề xuất. Thứ tự trình bày bám theo trình tự của chính quy trình: từ mạng
phân loại (\cref{sec:cnn}), qua kiến trúc phân đoạn (\cref{sec:unet}) và hàm mất
mát (\cref{sec:loss}), đến các phương pháp giải thích (\cref{sec:xai}), khung học
giám sát yếu (\cref{sec:weaksup}), mô hình nền tảng \SAM{} (\cref{sec:sam}), các
kỹ thuật xử lý ảnh cổ điển (\cref{sec:xulyanh}) và cuối cùng là hệ chỉ số đánh
giá (\cref{sec:metrics}).

\section{Mạng nơ-ron tích chập và học chuyển giao}
\label{sec:cnn}

\subsection{Mạng nơ-ron tích chập}

Mạng nơ-ron tích chập (\en{Convolutional Neural Network}, CNN) là kiến trúc chủ
đạo trong thị giác máy tính hiện đại. Khác với mạng kết nối đầy đủ, CNN khai thác
hai giả định cấu trúc phù hợp với dữ liệu ảnh: \emph{tính cục bộ} --- một điểm
ảnh liên quan mật thiết nhất với các điểm lân cận --- và \emph{tính bất biến tịnh
tiến} --- một đặc trưng như cạnh hay góc mang cùng ý nghĩa dù xuất hiện ở vị trí
nào trong ảnh.

Phép tích chập hai chiều giữa ảnh đầu vào $X$ và nhân tích chập $K$ kích thước
$m \times n$ được định nghĩa:

\begin{equation}
    (X * K)(i, j) = \sum_{a=0}^{m-1} \sum_{b=0}^{n-1} X(i + a,\, j + b) \cdot K(a, b)
    \label{eq:convolution}
\end{equation}

Trọng số của nhân $K$ được chia sẻ trên toàn bộ ảnh, làm giảm mạnh số tham số so
với mạng kết nối đầy đủ và đồng thời áp đặt tính bất biến tịnh tiến vào chính kiến
trúc. Một khối tích chập điển hình gồm phép tích chập, chuẩn hoá theo lô
(\en{batch normalization}) và hàm kích hoạt phi tuyến. Chồng nhiều khối như vậy
tạo ra một phân cấp đặc trưng: các tầng nông học cạnh và kết cấu, các tầng sâu học
các bộ phận và khái niệm ngữ nghĩa. Trong bài toán của luận văn, phân cấp này
tương ứng với việc mạng học từ kết cấu bề mặt lá ở tầng nông đến hình thái đốm
bệnh ở tầng sâu.

\subsection{Học chuyển giao}

Bộ dữ liệu của luận văn có \num{3104} ảnh huấn luyện --- một con số nhỏ so với
quy mô cần thiết để huấn luyện CNN từ đầu. Huấn luyện từ trạng thái khởi tạo ngẫu
nhiên trên tập nhỏ như vậy hầu như chắc chắn dẫn tới quá khớp (\en{overfitting}):
mạng ghi nhớ tập huấn luyện thay vì học quy luật khái quát.

Học chuyển giao (\en{transfer learning}) giải quyết vấn đề này bằng cách khởi tạo
mạng từ trọng số đã được huấn luyện trên một tập dữ liệu lớn, phổ biến nhất là
ImageNet với hơn 1,2 triệu ảnh thuộc 1000 lớp~\cite{deng2009imagenet}. Cơ sở của
cách làm là các đặc trưng ở tầng nông --- cạnh, góc, kết cấu, chuyển sắc màu ---
mang tính tổng quát và hữu ích cho hầu hết bài toán thị giác, bất kể miền dữ liệu.

Có hai chiến lược chính:

\begin{itemize}
    \item \textbf{Trích xuất đặc trưng} (\en{feature extraction}): đóng băng toàn
          bộ phần trích đặc trưng, chỉ huấn luyện tầng phân loại cuối. Phù hợp khi
          dữ liệu đích rất nhỏ hoặc rất giống miền nguồn.
    \item \textbf{Tinh chỉnh toàn bộ} (\en{full fine-tuning}): cập nhật mọi trọng
          số với tốc độ học nhỏ. Phù hợp khi dữ liệu đích đủ lớn và khác biệt đáng
          kể so với miền nguồn.
\end{itemize}

Ảnh lá bệnh khác xa phân bố ảnh ImageNet về cả nội dung lẫn kết cấu, trong khi
\num{3104} ảnh là đủ để tránh quên thảm hoạ (\en{catastrophic forgetting}) nếu
dùng tốc độ học đủ nhỏ. Vì vậy luận văn chọn tinh chỉnh toàn bộ với tốc độ học
$10^{-4}$; lý do lựa chọn được biện luận chi tiết ở \cref{chap:phuongphap}.

\subsection{Kiến trúc EfficientNet}

Trước EfficientNet, việc mở rộng CNN để tăng độ chính xác thường được thực hiện
theo một chiều đơn lẻ: tăng độ sâu (số tầng), tăng độ rộng (số kênh), hoặc tăng độ
phân giải ảnh đầu vào. Tan và Le~\cite{tan2019efficientnet} chỉ ra rằng ba chiều
này không độc lập --- mạng sâu hơn cần độ phân giải cao hơn để tận dụng vùng tiếp
nhận lớn, và cần nhiều kênh hơn để biểu diễn các mẫu tinh vi hơn --- và đề xuất
\emph{mở rộng hợp thành} (\en{compound scaling}), mở rộng đồng thời ba chiều theo
một hệ số duy nhất $\phi$:

\begin{equation}
    \text{độ sâu: } d = \alpha^{\phi}, \qquad
    \text{độ rộng: } w = \beta^{\phi}, \qquad
    \text{độ phân giải: } r = \gamma^{\phi}
    \label{eq:compound_scaling}
\end{equation}

với ràng buộc $\alpha \cdot \beta^{2} \cdot \gamma^{2} \approx 2$ và
$\alpha, \beta, \gamma \geq 1$. Ràng buộc này bảo đảm chi phí tính toán tăng xấp
xỉ $2^{\phi}$ lần khi hệ số hợp thành tăng $\phi$ đơn vị --- một cách kiểm soát
ngân sách tính toán rõ ràng. Các hệ số $\alpha, \beta, \gamma$ được xác định bằng
tìm kiếm lưới trên mạng cơ sở.

Mạng cơ sở EfficientNet-B0 được tìm ra bằng tìm kiếm kiến trúc tự động
(\en{neural architecture search}), với khối xây dựng chính là MBConv --- khối
tích chập nghịch đảo có nút cổ chai (\en{inverted residual bottleneck}) kết hợp
cơ chế \en{squeeze-and-excitation} để hiệu chỉnh trọng số giữa các kênh. Nhờ tích
chập tách theo chiều sâu (\en{depthwise separable convolution}), EfficientNet-B0
chỉ có khoảng \num{5.3} triệu tham số nhưng đạt độ chính xác cao hơn nhiều mạng
lớn hơn nó vài lần, chẳng hạn ResNet-50~\cite{he2016resnet} với khoảng
\num{25} triệu tham số.

Trong luận văn, EfficientNet-B0 giữ vai trò kép, và đây là một lựa chọn thiết kế
có chủ đích chứ không phải sự trùng hợp:

\begin{enumerate}
    \item Làm \textbf{mạng phân loại} ở Giai đoạn 1. Sau khi thay tầng phân loại
          1000 lớp bằng tầng 5 lớp, mô hình có \num{4013953} tham số.
    \item Làm \textbf{bộ mã hoá} (\en{encoder}) cho mạng phân đoạn ở Giai đoạn 4,
          với phần trích đặc trưng gồm \num{4007548} tham số.
\end{enumerate}

Việc dùng chung một kiến trúc xương sống cho cả hai giai đoạn bảo đảm bản đồ chú
ý sinh ra ở Giai đoạn 2 và đặc trưng mà mạng phân đoạn khai thác ở Giai đoạn 4
đến từ cùng một không gian biểu diễn, giúp quy trình nhất quán về mặt khái niệm.

\section{Kiến trúc phân đoạn ảnh}
\label{sec:unet}

\subsection{Từ phân loại đến phân đoạn}

Phân loại ảnh gán một nhãn cho toàn ảnh; phân đoạn ngữ nghĩa
(\en{semantic segmentation}) gán một nhãn cho \emph{từng điểm ảnh}. Chuyển từ bài
toán thứ nhất sang bài toán thứ hai đặt ra một mâu thuẫn cốt lõi: mạng phân loại
liên tục giảm độ phân giải không gian để mở rộng vùng tiếp nhận và trích xuất ngữ
nghĩa, nhưng phân đoạn lại đòi hỏi đầu ra có độ phân giải bằng đúng ảnh gốc. Với
đầu vào $224 \times 224$, bản đồ đặc trưng ở tầng cuối của EfficientNet-B0 chỉ
còn $7 \times 7$ --- mất mát thông tin không gian tới 1024 lần về diện tích.

\subsection{U-Net}

U-Net~\cite{ronneberger2015unet} giải quyết mâu thuẫn trên bằng kiến trúc mã
hoá--giải mã đối xứng hình chữ U. Nhánh mã hoá giảm dần độ phân giải để thu ngữ
nghĩa; nhánh giải mã tăng dần độ phân giải để khôi phục chi tiết không gian. Đóng
góp quyết định là các \emph{kết nối tắt} (\en{skip connection}) nối trực tiếp mỗi
tầng mã hoá với tầng giải mã tương ứng:

\begin{equation}
    x^{i}_{\text{dec}} = \mathcal{F}\big(\,[\,\mathcal{U}(x^{i+1}_{\text{dec}}) \;;\; x^{i}_{\text{enc}}\,]\,\big)
    \label{eq:unet_skip}
\end{equation}

trong đó $\mathcal{U}(\cdot)$ là phép tăng mẫu, $[\,\cdot\,;\,\cdot\,]$ là phép
nối theo chiều kênh và $\mathcal{F}(\cdot)$ là khối tích chập. Nhờ kết nối tắt,
thông tin không gian độ phân giải cao ở nhánh mã hoá được đưa thẳng sang nhánh
giải mã mà không phải đi qua nút cổ chai, cho phép khôi phục biên chính xác. U-Net
ban đầu được thiết kế cho ảnh y sinh --- bài toán cũng có đặc điểm dữ liệu huấn
luyện ít, tương tự bối cảnh của luận văn.

\subsection{UNet++}

Zhou và cộng sự~\cite{zhou2018unetpp} chỉ ra một điểm yếu của U-Net: kết nối tắt
nối trực tiếp đặc trưng \emph{nông} của bộ mã hoá với đặc trưng \emph{sâu} của bộ
giải mã, tạo ra \emph{khoảng cách ngữ nghĩa} (\en{semantic gap}) giữa hai luồng
thông tin có mức trừu tượng rất khác nhau. Việc nối trực tiếp buộc bộ giải mã phải
tự dung hoà hai nguồn không đồng nhất.

\UNetpp{} thu hẹp khoảng cách này bằng các kết nối tắt \emph{lồng nhau, dày đặc}:
thay vì một đường nối thẳng, đặc trưng của bộ mã hoá đi qua một chuỗi khối tích
chập trung gian trước khi tới bộ giải mã. Node $x^{i,j}$ ($i$ là chỉ số tầng, $j$
là chỉ số dọc theo đường nối) được tính:

\begin{equation}
    x^{i,j} =
    \begin{cases}
        \mathcal{F}\big(x^{i-1,j}\big) & \text{nếu } j = 0 \\[6pt]
        \mathcal{F}\Big(\big[\,[x^{i,k}]_{k=0}^{j-1},\; \mathcal{U}(x^{i+1,j-1})\,\big]\Big) & \text{nếu } j > 0
    \end{cases}
    \label{eq:unetpp}
\end{equation}

Mỗi node nhận toàn bộ các node trước đó cùng tầng cộng với đặc trưng đã tăng mẫu
từ tầng dưới. Kết quả là đặc trưng được tinh chỉnh dần qua nhiều bước trung gian
thay vì nhảy một bước, và bộ giải mã nhận được biểu diễn ở nhiều mức trừu tượng
khác nhau.

Đặc điểm này đặc biệt phù hợp với bài toán của luận văn. Bệnh đốm lá tạo ra các
tổn thương \emph{nhiều kích thước khác nhau trong cùng một ảnh}: đốm nhỏ vài chục
điểm ảnh nằm cạnh mảng hoại tử lớn. Kiến trúc lồng nhau cho phép mô hình kết hợp
thông tin đa tỉ lệ tự nhiên hơn so với kết nối tắt đơn.

Chi phí của \UNetpp{} là khiêm tốn. Với cùng bộ mã hoá EfficientNet-B0, mô hình
\UNetpp{} có \num{6569581} tham số so với \num{6251469} tham số của U-Net thường
--- chỉ tăng \num{318112} tham số, tương đương 5,1\%. Trong đó bộ mã hoá chiếm
\num{4007548} tham số và bộ giải mã chiếm \num{2561888} tham số. Mức tăng nhỏ này
là lý do luận văn chọn \UNetpp{} mà không lo ngại về ngân sách tính toán.

\section{Hàm mất mát cho dữ liệu mất cân bằng}
\label{sec:loss}

\subsection{Bản chất của mất cân bằng trong phân đoạn tổn thương}

Trong bài toán phân đoạn nhị phân vùng bệnh, số điểm ảnh thuộc lớp nền lớn hơn
nhiều lần số điểm ảnh thuộc lớp tổn thương. Với độ phủ mục tiêu 18--22\% như
thiết kế ở \cref{chap:phuongphap}, tỉ lệ nền trên tổn thương xấp xỉ 4:1; với các
ảnh bệnh nhẹ, tỉ lệ này còn cao hơn nhiều. Mất cân bằng dẫn tới một nghiệm suy
biến nguy hiểm: mô hình dự đoán toàn bộ ảnh là nền vẫn đạt độ chính xác theo điểm
ảnh rất cao nhưng hoàn toàn vô dụng.

\subsection{Binary Cross-Entropy}

Hàm mất mát cơ bản cho phân loại nhị phân theo từng điểm ảnh:

\begin{equation}
    \mathcal{L}_{\text{BCE}} = -\frac{1}{N}\sum_{i=1}^{N}
        \Big[ y_i \log(p_i) + (1 - y_i)\log(1 - p_i) \Big]
    \label{eq:bce}
\end{equation}

trong đó $y_i \in \{0,1\}$ là nhãn và $p_i \in (0,1)$ là xác suất dự đoán của
điểm ảnh thứ $i$. Hạn chế của BCE là xử lý mọi điểm ảnh với trọng số như nhau:
khi lớp nền áp đảo, tổng đóng góp của các điểm nền lấn át tín hiệu từ vùng tổn
thương.

\subsection{Dice Loss}

Milletari và cộng sự~\cite{milletari2016vnet} đề xuất tối ưu trực tiếp hệ số Dice
--- một đại lượng đo độ chồng lấp, không phụ thuộc vào số điểm ảnh nền:

\begin{equation}
    \mathcal{L}_{\text{Dice}} = 1 - \frac{2\sum_{i} p_i y_i + \epsilon}
                                         {\sum_{i} p_i + \sum_{i} y_i + \epsilon}
    \label{eq:dice}
\end{equation}

Hằng số $\epsilon$ giữ cho biểu thức xác định khi cả dự đoán lẫn nhãn đều rỗng
--- trường hợp thực sự xảy ra trong luận văn, vì \num{683} ảnh lá khoẻ được gán
mặt nạ toàn số không (\cref{chap:phuongphap}). Ưu điểm của Dice là tối ưu trực
tiếp đại lượng dùng để đánh giá; nhược điểm là gradient kém ổn định khi vùng đích
rất nhỏ.

\subsection{Focal Loss}

Lin và cộng sự~\cite{lin2017focal} tiếp cận mất cân bằng theo hướng khác: thay vì
đổi đại lượng tối ưu, họ \emph{giảm trọng số} của các mẫu đã được phân loại tốt:

\begin{equation}
    \mathcal{L}_{\text{Focal}} = -\alpha \,(1 - p_t)^{\gamma} \log(p_t),
    \qquad
    p_t = \begin{cases} p & \text{nếu } y = 1 \\ 1 - p & \text{nếu } y = 0 \end{cases}
    \label{eq:focal}
\end{equation}

Hệ số điều biến $(1-p_t)^{\gamma}$ là mấu chốt. Với một điểm ảnh nền dễ, mô hình
dự đoán $p_t = 0{,}95$, hệ số bằng $0{,}05^{2} = 0{,}0025$ khi $\gamma = 2$ ---
đóng góp của nó bị giảm 400 lần. Ngược lại, một điểm ảnh khó với $p_t = 0{,}5$
chỉ bị giảm 4 lần. Nhờ vậy gradient tập trung vào vùng biên tổn thương --- chính
là nơi mô hình cần học nhất. Tham số $\alpha$ cân bằng thêm giữa hai lớp. Luận văn
dùng $\alpha = 0{,}25$ và $\gamma = 2{,}0$ theo giá trị khuyến nghị của tác giả
gốc.

\subsection{Hàm mất mát kết hợp FocalDice}

Hai hàm trên bổ khuyết cho nhau: Focal Loss hoạt động ở mức \emph{từng điểm ảnh},
điều tiết đóng góp theo độ khó; Dice Loss hoạt động ở mức \emph{toàn vùng}, tối
ưu độ chồng lấp tổng thể. Luận văn sử dụng tổ hợp:

\begin{equation}
    \mathcal{L}_{\text{FocalDice}} = \mathcal{L}_{\text{Focal}} + \mathcal{L}_{\text{Dice}}
    \label{eq:focaldice}
\end{equation}

Cần lưu ý một hệ quả thực nghiệm quan trọng: do thành phần Focal chủ động thu nhỏ
gradient của phần lớn điểm ảnh, giá trị mất mát tổng thể giảm chậm hơn nhiều so
với BCE + Dice. Điều này khiến \emph{tốc độ giảm của hàm mất mát không còn là chỉ
báo đáng tin cậy cho tiến bộ của mô hình} --- một quan sát dẫn tới hệ quả trực
tiếp về tiêu chí dừng sớm, được phân tích ở \cref{chap:thucnghiem}.

\subsection{Thuật toán tối ưu}

Luận văn sử dụng hai bộ tối ưu. Adam~\cite{kingma2015adam} ước lượng thích nghi
mô-men bậc một và bậc hai của gradient, cho tốc độ hội tụ nhanh mà ít cần điều
chỉnh tốc độ học. AdamW~\cite{loshchilov2019adamw} sửa một khiếm khuyết của Adam:
trong Adam, suy giảm trọng số (\en{weight decay}) được cộng vào gradient nên bị
chia cho ước lượng mô-men bậc hai, làm cường độ chính quy hoá thay đổi không kiểm
soát giữa các tham số. AdamW tách suy giảm trọng số khỏi bước cập nhật gradient,
khôi phục đúng ý nghĩa chính quy hoá.

Về lịch tốc độ học, luận văn dùng hai chiến lược khác nhau cho hai nhóm thực
nghiệm. \en{ReduceLROnPlateau} giảm tốc độ học khi chỉ số theo dõi ngừng cải
thiện --- đơn giản và không cần biết trước tổng số bước. \en{CosineAnnealingLR},
xuất phát từ công trình của Loshchilov và Hutter~\cite{loshchilov2017sgdr}, giảm
tốc độ học theo quy luật cô-sin:

\begin{equation}
    \eta_t = \eta_{\min} + \tfrac{1}{2}\big(\eta_{\max} - \eta_{\min}\big)
             \left(1 + \cos\!\Big(\frac{T_{\text{cur}}}{T_{\max}}\pi\Big)\right)
    \label{eq:cosine}
\end{equation}

Lịch cô-sin giảm chậm ở đầu để mô hình khám phá không gian nghiệm, rồi giảm nhanh
về cuối để tinh chỉnh --- và quan trọng là không bao giờ bị kẹt như
\en{ReduceLROnPlateau} khi chỉ số theo dõi dao động.

\section{Trí tuệ nhân tạo khả diễn giải}
\label{sec:xai}

\subsection{Class Activation Map}

Zhou và cộng sự~\cite{zhou2016cam} đề xuất CAM, phương pháp đầu tiên định vị được
vùng ảnh mà mạng phân loại dựa vào để ra quyết định. Với mạng kết thúc bằng gộp
trung bình toàn cục (\en{global average pooling}) rồi tới một tầng tuyến tính,
điểm số của lớp $c$ là $y^{c} = \sum_{k} w^{c}_{k} \cdot \frac{1}{Z}\sum_{i,j} A^{k}_{ij}$.
Hoán đổi thứ tự hai phép tổng cho ra bản đồ định vị:

\begin{equation}
    M^{c}(i, j) = \sum_{k} w^{c}_{k} \, A^{k}_{ij}
    \label{eq:cam}
\end{equation}

Hạn chế nghiêm trọng của CAM là nó \emph{đòi hỏi kiến trúc cụ thể}: mạng bắt buộc
phải có gộp trung bình toàn cục ngay trước tầng phân loại. Mạng nào không thoả
mãn phải sửa kiến trúc rồi huấn luyện lại.

\subsection{Grad-CAM}

Selvaraju và cộng sự~\cite{selvaraju2017gradcam} tổng quát hoá CAM bằng nhận xét
rằng trọng số $w^{c}_{k}$ trong \cref{eq:cam} có thể được thay bằng gradient trung
bình --- một đại lượng tính được ở \emph{bất kỳ} tầng tích chập nào của
\emph{bất kỳ} kiến trúc nào, không cần huấn luyện lại:

\begin{equation}
    \alpha^{c}_{k} = \frac{1}{Z}\sum_{i}\sum_{j} \frac{\partial y^{c}}{\partial A^{k}_{ij}}
    \label{eq:gradcam_alpha}
\end{equation}

\begin{equation}
    L^{c}_{\GradCAM} = \operatorname{ReLU}\!\left(\sum_{k} \alpha^{c}_{k} A^{k}\right)
    \label{eq:gradcam}
\end{equation}

Hàm ReLU ở \cref{eq:gradcam} giữ lại chỉ những vùng có ảnh hưởng \emph{dương} tới
điểm số của lớp quan tâm; vùng có ảnh hưởng âm là bằng chứng chống lại lớp đó nên
bị loại. Tính tổng quát về kiến trúc khiến \GradCAM{} trở thành phương pháp giải
thích được dùng rộng rãi nhất cho CNN.

\subsection{Grad-CAM++}

\GradCAM{} bộc lộ điểm yếu khi một ảnh chứa \emph{nhiều thể hiện} của cùng một
lớp. Do \cref{eq:gradcam_alpha} lấy trung bình gradient trên toàn bản đồ đặc
trưng, một vùng có gradient lớn sẽ chi phối giá trị trung bình và làm các vùng
khác bị san phẳng. Hệ quả là bản đồ nhiệt có xu hướng chỉ làm nổi bật thể hiện rõ
nhất và bỏ sót các thể hiện còn lại.

Đây chính xác là tình huống của bệnh đốm lá: một chiếc lá thường mang hàng chục
đốm bệnh phân tán, mỗi đốm là một thể hiện riêng của cùng lớp bệnh.

Chattopadhay và cộng sự~\cite{chattopadhay2018gradcampp} khắc phục bằng cách thay
trọng số trung bình đồng đều bằng trọng số riêng cho từng vị trí không gian, dựa
trên đạo hàm bậc hai và bậc ba:

\begin{equation}
    \alpha^{kc}_{ij} = \frac{\dfrac{\partial^{2} y^{c}}{\partial (A^{k}_{ij})^{2}}}
    {2\dfrac{\partial^{2} y^{c}}{\partial (A^{k}_{ij})^{2}}
     + \displaystyle\sum_{a}\sum_{b} A^{k}_{ab} \dfrac{\partial^{3} y^{c}}{\partial (A^{k}_{ij})^{3}}}
    \label{eq:gradcampp_alpha}
\end{equation}

\begin{equation}
    w^{c}_{k} = \sum_{i}\sum_{j} \alpha^{kc}_{ij} \cdot
                \operatorname{ReLU}\!\left(\frac{\partial y^{c}}{\partial A^{k}_{ij}}\right)
    \label{eq:gradcampp_w}
\end{equation}

Mỗi vị trí $(i,j)$ nay có trọng số riêng $\alpha^{kc}_{ij}$, nên nhiều vùng kích
hoạt cùng tồn tại trong bản đồ kết quả thay vì bị một vùng lấn át. Đó là lý do
luận văn chuyển từ \GradCAM{} sang \GradCAMpp{} khi nâng cấp nhãn giả từ V1 lên
V2.

\subsection{Hợp nhất đa tỉ lệ}

Bản đồ kích hoạt kế thừa độ phân giải của tầng được chọn, kéo theo một sự đánh
đổi trực tiếp:

\begin{itemize}
    \item Tầng \textbf{sâu} (với đầu vào $224 \times 224$, EfficientNet-B0 cho
          bản đồ $7 \times 7$): giàu ngữ nghĩa, biết chắc ``đây là vùng bệnh'',
          nhưng khi phóng lên $224 \times 224$ thì mỗi ô tương ứng một khối
          $32 \times 32$ điểm ảnh --- biên trở nên rất thô.
    \item Tầng \textbf{nông hơn} ($14 \times 14$): định vị chi tiết hơn gấp bốn
          lần về diện tích, nhưng ngữ nghĩa yếu hơn nên dễ kích hoạt nhầm trên
          kết cấu nền.
\end{itemize}

Cách dung hoà là hợp nhất tuyến tính bản đồ từ hai tầng:

\begin{equation}
    H_{\text{final}} = \lambda \, H_{\text{fine}} + (1 - \lambda) \, H_{\text{coarse}}
    \label{eq:multiscale}
\end{equation}

Luận văn dùng $\lambda = 0{,}6$, thiên về tầng tinh để mặt nạ sắc nét hơn nhưng
vẫn giữ đủ ngữ cảnh ngữ nghĩa từ tầng thô. Lựa chọn cụ thể của các tầng được trình
bày ở \cref{chap:phuongphap}.

\section{Học giám sát yếu và gán nhãn giả}
\label{sec:weaksup}

\subsection{Phân loại các mức giám sát}

Theo mức độ chi tiết của nhãn có sẵn, các bài toán phân đoạn được chia thành:

\begin{itemize}
    \item \textbf{Giám sát đầy đủ}: mỗi ảnh có mặt nạ ở mức điểm ảnh. Cho kết quả
          tốt nhất nhưng chi phí gán nhãn cao nhất.
    \item \textbf{Giám sát yếu} (\en{weakly-supervised}): chỉ có nhãn thô hơn ---
          nhãn mức ảnh, hộp bao, nét vẽ nguệch ngoạc hoặc điểm đánh dấu.
    \item \textbf{Giám sát bán phần} (\en{semi-supervised}): một phần nhỏ có mặt
          nạ đầy đủ, phần lớn không có nhãn.
    \item \textbf{Không giám sát}: hoàn toàn không có nhãn.
\end{itemize}

Luận văn thuộc nhóm thứ hai, ở dạng khó nhất: chỉ có nhãn \emph{mức ảnh}.

\subsection{Gán nhãn giả}

Gán nhãn giả (\en{pseudo-labeling}), được Lee~\cite{lee2013pseudolabel} giới thiệu
cho học bán giám sát, là kỹ thuật dùng chính dự đoán của mô hình làm nhãn huấn
luyện cho vòng học tiếp theo. Trong bối cảnh chuyển từ phân loại sang phân đoạn,
ý tưởng được vận dụng như sau: bản đồ kích hoạt của mạng phân loại được nhị phân
hoá thành mặt nạ, và mặt nạ đó đóng vai trò nhãn giả để huấn luyện mạng phân đoạn.

Chuỗi suy luận có thể tóm tắt: \emph{nhãn mức ảnh $\rightarrow$ mạng phân loại
$\rightarrow$ bản đồ kích hoạt $\rightarrow$ ngưỡng hoá $\rightarrow$ mặt nạ giả
$\rightarrow$ mạng phân đoạn}. Barbedo~\cite{barbedo2019lesions} đã chứng minh
tính khả thi của hướng tiếp cận dựa trên tổn thương này trong lĩnh vực bệnh cây
trồng.

\subsection{Hạn chế cố hữu cần ghi nhận}

Khung phương pháp trên mang một hạn chế bản chất, không thể khắc phục bằng kỹ
thuật, và luận văn coi việc nêu rõ nó là một phần của đóng góp:

\begin{quote}
Bản đồ kích hoạt trả lời câu hỏi \emph{``mô hình nhìn vào đâu để phân loại''}, chứ
không trả lời câu hỏi \emph{``vùng bệnh nằm ở đâu''}. Hai câu hỏi này tương quan
với nhau nhưng không đồng nhất.
\end{quote}

Mạng phân loại chỉ cần \emph{đủ} bằng chứng để phân biệt các lớp, nên nó có thể
tập trung vào một vài đốm điển hình nhất mà bỏ qua phần còn lại --- vẫn phân loại
đúng. Vì vậy nhãn giả sinh ra từ bản đồ kích hoạt thường chưa phủ hết vùng bệnh
thực. Hệ quả kéo theo là mọi chỉ số đánh giá tính trên nhãn giả đều đo mức tái
tạo nhãn giả chứ không đo độ chính xác bệnh học --- vấn đề được phân tích trọn vẹn
ở \cref{sec:caveat}.

\section{Mô hình nền tảng Segment Anything}
\label{sec:sam}

\subsection{Kiến trúc và khả năng}

\SAM{}~\cite{kirillov2023sam} là mô hình nền tảng (\en{foundation model}) cho bài
toán phân đoạn, được Meta AI huấn luyện trên bộ dữ liệu SA-1B gồm khoảng 11 triệu
ảnh và hơn 1 tỉ mặt nạ. Kiến trúc gồm ba thành phần:

\begin{itemize}
    \item \textbf{Bộ mã hoá ảnh}: một Vision Transformer, chạy một lần cho mỗi
          ảnh và cho ra biểu diễn nhúng dùng lại được.
    \item \textbf{Bộ mã hoá gợi ý}: mã hoá gợi ý đầu vào dưới dạng điểm, hộp bao,
          mặt nạ thô hoặc văn bản.
    \item \textbf{Bộ giải mã mặt nạ}: kết hợp hai nguồn nhúng trên để sinh mặt nạ,
          rất nhẹ nên cho phép thử nhiều gợi ý khác nhau với chi phí thấp.
\end{itemize}

Luận văn dùng biến thể ViT-B --- bản nhỏ nhất trong ba biến thể ViT-B/ViT-L/ViT-H,
với tệp trọng số khoảng 375 MB.

\subsection{Vai trò bộ tinh chỉnh nhãn}

Điểm mạnh của \SAM{} là khả năng khái quát không cần huấn luyện lại
(\en{zero-shot}) và chất lượng biên rất cao. Điều này gợi ý một cách dùng khác với
mục đích ban đầu: thay vì dùng \SAM{} để phân đoạn trực tiếp, ta dùng mặt nạ giả
thô làm gợi ý và để \SAM{} tinh chỉnh lại biên. Về nguyên tắc, cách làm này kết
hợp được ưu điểm của cả hai nguồn --- tri thức về \emph{vị trí} bệnh từ bản đồ
kích hoạt, và tri thức về \emph{biên} vật thể từ \SAM{}.

Tuy nhiên, cách dùng này chứa một rủi ro then chốt. \SAM{} được huấn luyện để phân
đoạn \emph{đơn vị ngữ nghĩa tự nhiên} của cảnh. Khi nhận gợi ý là vài điểm nằm
trên một đốm bệnh, đơn vị ngữ nghĩa mà \SAM{} nhận diện có thể không phải là đốm
bệnh mà là \emph{cả chiếc lá} chứa đốm đó. Cơ chế bảo vệ dựa trên IoU thường được
dùng để phát hiện sai lệch, nhưng như phân tích ở \cref{chap:thucnghiem}, cơ chế
này chỉ chặn được trường hợp mặt nạ \emph{trôi} sang vùng khác chứ không chặn được
trường hợp mặt nạ \emph{nở} ra bao trọn mặt nạ ban đầu. Đây là một trong những
phát hiện thực nghiệm chính của luận văn.

\section{Kỹ thuật xử lý ảnh cổ điển}
\label{sec:xulyanh}

Bản đồ kích hoạt là ảnh xám liên tục, trong khi nhãn phân đoạn phải là mặt nạ nhị
phân sạch. Khoảng cách giữa hai dạng biểu diễn được lấp bằng các kỹ thuật xử lý
ảnh cổ điển dưới đây.

\subsection{Ngưỡng hoá Otsu}

Phương pháp Otsu~\cite{otsu1979threshold} chọn ngưỡng tự động bằng cách cực đại
hoá phương sai giữa hai lớp:

\begin{equation}
    \sigma_b^{2}(t) = \omega_0(t)\,\omega_1(t)\,\big[\mu_0(t) - \mu_1(t)\big]^{2}
    \label{eq:otsu}
\end{equation}

với $\omega_0, \omega_1$ là tỉ lệ điểm ảnh và $\mu_0, \mu_1$ là giá trị trung bình
của hai lớp khi cắt tại ngưỡng $t$. Otsu hoạt động tốt khi biểu đồ mức xám có
\emph{hai đỉnh} tách bạch. Giả định này bị vi phạm khi bản đồ nhiệt được hợp nhất
đa tỉ lệ theo \cref{eq:multiscale}: phép hợp nhất làm phân bố trở nên trơn và đơn
đỉnh, khiến Otsu chọn ngưỡng quá thấp và gây phân đoạn thừa. Đây là lý do trực
tiếp khiến luận văn thay Otsu bằng ngưỡng phân vị theo lớp, trình bày ở
\cref{chap:phuongphap}.

\subsection{Phép toán hình thái}

Bốn phép toán hình thái cơ bản trên ảnh nhị phân $A$ với phần tử cấu trúc $B$:

\begin{align}
    \text{Co (\en{erosion}):} \quad & A \ominus B = \{z \mid B_z \subseteq A\} \label{eq:erosion}\\
    \text{Giãn (\en{dilation}):} \quad & A \oplus B = \{z \mid B_z \cap A \neq \emptyset\} \label{eq:dilation}\\
    \text{Mở (\en{opening}):} \quad & A \circ B = (A \ominus B) \oplus B \label{eq:opening}\\
    \text{Đóng (\en{closing}):} \quad & A \bullet B = (A \oplus B) \ominus B \label{eq:closing}
\end{align}

Trong quy trình của luận văn, phép \emph{đóng} lấp các lỗ nhỏ bên trong vùng bệnh
--- cần thiết vì đốm bệnh thường có tâm sáng hơn viền nên bị ngưỡng hoá cắt rỗng
--- còn phép \emph{mở} loại các điểm nhiễu rải rác bên ngoài vùng chính. Hình dạng
phần tử cấu trúc có ý nghĩa thực tế: phần tử \emph{elip} bám sát hình dạng tròn
hoặc bầu dục của đốm bệnh tốt hơn phần tử \emph{chữ nhật}, và đây là một trong các
cải tiến từ nhãn giả V1 sang V2.

\subsection{Lọc thành phần liên thông}

Sau ngưỡng hoá, mặt nạ thường gồm nhiều vùng liên thông rời rạc. Phân tích thành
phần liên thông (\en{connected component analysis}) gán nhãn cho từng vùng và cho
phép lọc theo diện tích hoặc theo thứ hạng. Chiến lược giữ lại $k$ vùng lớn nhất
có diện tích vượt một ngưỡng tối thiểu phù hợp với đặc thù bệnh đốm lá: cần giữ
\emph{nhiều} vùng vì bệnh biểu hiện thành nhiều đốm, nhưng phải loại các mảnh vụn
nhiễu mà phép mở hình thái không xử lý hết.

\subsection{Lọc song phương}

Lọc song phương (\en{bilateral filter}) của Tomasi và
Manduchi~\cite{tomasi1998bilateral} làm trơn ảnh nhưng bảo toàn biên, nhờ nhân
trọng số tính trên cả khoảng cách không gian lẫn chênh lệch cường độ:

\begin{equation}
    I^{\text{filt}}(p) = \frac{1}{W_p}\sum_{q \in \Omega} I(q)\,
        \underbrace{G_{\sigma_s}\big(\lVert p - q \rVert\big)}_{\text{gần về không gian}}\;
        \underbrace{G_{\sigma_r}\big(|I(p) - I(q)|\big)}_{\text{giống về cường độ}}
    \label{eq:bilateral}
\end{equation}

Hai điểm ảnh chỉ ảnh hưởng lẫn nhau khi vừa gần nhau về vị trí vừa tương đồng về
cường độ. Áp dụng lên bản đồ nhiệt, bộ lọc này khiến kích hoạt ``lan'' theo các
vùng đồng màu và dừng lại tại ranh giới màu sắc thực của tổn thương.

\subsection{Trường ngẫu nhiên có điều kiện}

Dense CRF~\cite{krahenbuhl2011densecrf} tinh chỉnh kết quả phân đoạn bằng cách
cực tiểu hoá hàm năng lượng gồm thế đơn (dựa trên dự đoán ban đầu) và thế cặp
(khuyến khích các điểm ảnh gần nhau và giống màu nhận cùng nhãn):

\begin{equation}
    E(\mathbf{x}) = \sum_{i} \psi_u(x_i) + \sum_{i < j} \psi_p(x_i, x_j)
    \label{eq:crf}
\end{equation}

Luận văn có khảo sát Dense CRF như một hướng tinh chỉnh thay thế cho \SAM{}, nhưng
không đưa vào quy trình cuối vì lý do triển khai được nêu ở \cref{chap:thucnghiem}.

\section{Chỉ số đánh giá}
\label{sec:metrics}

\subsection{Chỉ số cho bài toán phân loại}

Từ ma trận nhầm lẫn, với TP, FP, FN lần lượt là số dương tính thật, dương tính giả
và âm tính giả:

\begin{equation}
    \text{Precision} = \frac{TP}{TP + FP}, \qquad
    \text{Recall} = \frac{TP}{TP + FN}, \qquad
    F_1 = \frac{2 \cdot P \cdot R}{P + R}
    \label{eq:prf}
\end{equation}

Với bài toán chẩn đoán bệnh cây trồng, \textbf{recall quan trọng hơn precision}:
bỏ sót một cây bệnh (âm tính giả) gây thiệt hại lan rộng, trong khi cảnh báo nhầm
một cây khoẻ (dương tính giả) chỉ tốn thêm một lần kiểm tra. Nhận định này định
hướng cách diễn giải kết quả ở \cref{chap:thucnghiem}.

Do bộ dữ liệu gần cân bằng (tỉ lệ mất cân bằng 1,33), luận văn báo cáo cả
\en{macro-average} --- trung bình không trọng số qua các lớp, phản ánh hiệu năng
trên lớp ít mẫu --- và \en{weighted-average}.

\subsection{Chỉ số cho bài toán phân đoạn}

Hai chỉ số chuẩn đo độ chồng lấp giữa mặt nạ dự đoán $X$ và mặt nạ tham chiếu $Y$:

\begin{equation}
    \text{IoU} = \frac{|X \cap Y|}{|X \cup Y|} = \frac{TP}{TP + FP + FN}
    \label{eq:iou}
\end{equation}

\begin{equation}
    \text{Dice} = \frac{2|X \cap Y|}{|X| + |Y|} = \frac{2\,TP}{2\,TP + FP + FN}
    \label{eq:dicecoef}
\end{equation}

Hai chỉ số liên hệ với nhau qua một song ánh đơn điệu:

\begin{equation}
    \text{Dice} = \frac{2\,\text{IoU}}{1 + \text{IoU}}, \qquad
    \text{IoU} = \frac{\text{Dice}}{2 - \text{Dice}}
    \label{eq:iou_dice_relation}
\end{equation}

Từ \cref{eq:iou_dice_relation} suy ra hai hệ quả cần lưu ý khi đọc kết quả. Thứ
nhất, Dice luôn lớn hơn hoặc bằng IoU, nên báo cáo Dice thay vì IoU luôn cho con
số ``đẹp'' hơn dù chất lượng không đổi. Thứ hai, do quan hệ là đơn điệu, việc xếp
hạng các mô hình theo IoU hay theo Dice luôn cho cùng thứ tự --- báo cáo cả hai
không cung cấp thêm thông tin về thứ hạng, mà chỉ giúp đối chiếu với các công bố
dùng chỉ số khác nhau.

Trong bài toán phân đoạn nhị phân, $F_1$ tính trên tập điểm ảnh trùng khớp hoàn
toàn với hệ số Dice --- điều này giải thích vì sao hai cột $F_1$ và Dice trong
các bảng kết quả ở \cref{chap:thucnghiem} luôn bằng nhau.

\subsection{Ghi chú về mốc tham chiếu}

Mọi công thức trong mục này đều giả định $Y$ là nhãn chuẩn (\en{ground truth}).
Trong luận văn, do không có nhãn thủ công ở mức điểm ảnh, $Y$ là \emph{nhãn giả}.
Các công thức vẫn tính được và vẫn có ý nghĩa so sánh nội bộ, nhưng ý nghĩa diễn
giải thay đổi hoàn toàn: chúng đo mức độ mô hình tái tạo nhãn giả, chứ không đo
độ chính xác so với tổn thương thực. \Cref{sec:caveat} phân tích trọn vẹn hệ quả
của sự thay đổi này.

\section{Nghiên cứu liên quan}
\label{sec:lienquan}

\subsection{Học sâu trong phát hiện bệnh cây trồng}

Mohanty và cộng sự~\cite{mohanty2016plantdisease} là công trình mở đường, đạt độ
chính xác 99,35\% trên bộ PlantVillage với 54.306 ảnh thuộc 38 lớp bệnh--cây
trồng. Tuy nhiên các tác giả cũng ghi nhận độ chính xác sụt giảm mạnh khi kiểm
thử trên ảnh chụp trong điều kiện thực địa --- do PlantVillage được chụp trên nền
đồng nhất trong phòng thí nghiệm. Đây là bài học nền tảng về khoảng cách giữa dữ
liệu chuẩn hoá và dữ liệu thực tế.

Ferentinos~\cite{ferentinos2018plantdisease} mở rộng khảo sát trên 87.848 ảnh với
25 loài cây, so sánh nhiều kiến trúc CNN và đạt độ chính xác cao nhất 99,53\%.
Kamilaris và Prenafeta-Boldú~\cite{kamilaris2018survey} tổng hợp 40 nghiên cứu
ứng dụng học sâu trong nông nghiệp và kết luận rằng học sâu vượt trội so với các
kỹ thuật xử lý ảnh truyền thống, đồng thời chỉ ra hai điểm nghẽn: thiếu dữ liệu
có nhãn chất lượng và thiếu khả năng diễn giải.

Fuentes và cộng sự~\cite{fuentes2017tomato} tiếp cận theo hướng phát hiện đối
tượng thay vì phân loại, cho phép định vị nhiều vùng bệnh trên cùng một ảnh cà
chua. Barbedo~\cite{barbedo2019lesions} lập luận rằng nhận dạng ở mức
\emph{từng tổn thương} thay vì mức toàn lá giúp tăng đáng kể lượng mẫu huấn luyện
hiệu dụng và cải thiện khả năng khái quát --- lập luận trực tiếp ủng hộ hướng
phân đoạn mà luận văn theo đuổi.

\subsection{Vị trí của luận văn}

Đối chiếu với các công trình trên, luận văn có ba điểm khác biệt:

\begin{enumerate}
    \item \textbf{Đối tượng}: sầu riêng là cây trồng chưa được nghiên cứu nhiều
          trong lĩnh vực phát hiện bệnh bằng học sâu, so với cà chua, lúa hay
          khoai tây.
    \item \textbf{Mức đầu ra}: hầu hết công trình dừng ở phân loại hoặc phát hiện
          bằng hộp bao; luận văn hướng tới phân đoạn ở mức điểm ảnh mà không dùng
          nhãn thủ công nào ở mức đó.
    \item \textbf{Thái độ với kết quả}: thay vì chỉ báo cáo chỉ số cao nhất đạt
          được, luận văn phân tích thẳng vào việc \emph{các chỉ số ấy có nghĩa gì
          và không có nghĩa gì} khi thiếu nhãn chuẩn --- một khía cạnh thường bị
          bỏ qua trong các nghiên cứu học giám sát yếu.
\end{enumerate}

\section{Kết luận chương}

Chương này đã trình bày nền tảng lý thuyết cho toàn bộ quy trình đề xuất:
EfficientNet với cơ chế mở rộng hợp thành làm mạng phân loại và bộ mã hoá;
\UNetpp{} với kết nối tắt lồng nhau làm mạng phân đoạn; tổ hợp Focal và Dice làm
hàm mất mát cho dữ liệu mất cân bằng; \GradCAMpp{} hợp nhất đa tỉ lệ để sinh bản
đồ kích hoạt; \SAM{} trong vai trò bộ tinh chỉnh biên; cùng các kỹ thuật xử lý ảnh
cổ điển để chuyển bản đồ nhiệt thành mặt nạ nhị phân.

Hai điểm lý thuyết sẽ được viện dẫn nhiều lần ở các chương sau, cần được ghi nhận
ngay: (i) bản đồ kích hoạt trả lời câu hỏi ``mô hình nhìn vào đâu'' chứ không phải
``bệnh nằm ở đâu'' (\cref{sec:weaksup}); và (ii) mọi chỉ số IoU/Dice đều được định
nghĩa \emph{tương đối} với một mốc tham chiếu, nên khi mốc tham chiếu thay đổi thì
các con số không còn so sánh trực tiếp được với nhau (\cref{sec:metrics}).

\Cref{chap:phuongphap} trình bày cách các thành phần lý thuyết trên được lắp ghép
thành một quy trình hoàn chỉnh, cùng lý do kỹ thuật của từng lựa chọn thiết kế.
